\chapter{Símbolos y parámetros de consulta}
\begin{longtable}{@{}P{31mm}P{82mm}P{28mm}@{}}
\caption{Nomenclatura del informe.}\\
\toprule\textbf{Símbolo} & \textbf{Significado} & \textbf{Unidad}\\\midrule\endfirsthead
\toprule\textbf{Símbolo} & \textbf{Significado} & \textbf{Unidad}\\\midrule\endhead
$\vct q$ & Ángulos geométricos del manipulador. & rad\\
$x,y,\psi$ & Pose planar de la base. & m, m, rad\\
$v_x,v_y,\omega$ & Velocidades expresadas en el cuerpo. & m/s, rad/s\\
$u_i,\omega_{w,i}$ & Velocidad periférica y angular de rueda. & m/s, rad/s\\
$r,D,L$ & Radio, diámetro y suma de semidistancias de base. & m\\
$N_i,\Delta n_i$ & Cuentas por vuelta e incremento de cuentas. & cuenta/vuelta; cuenta\\
$h,L_1,L_2,d$ & Longitudes equivalentes de la cadena del brazo. & m\\
$\mat T_{AB}$ & Transformación que lleva coordenadas de B a A. & Mixta\\
$\mat R_{AB},\vct t_{AB}$ & Rotación y traslación entre marcos. & Adimensional; m\\
$W,B,M,E$ & Mundo, base, montura y herramienta. & Marcos\\
$\mat J$ & Jacobiano geométrico, orden lineal--angular. & Según filas\\
$\vct F,\vct M$ & Fuerza y momento externos. & N, N\textperiodcentered{}m\\
$\vct\tau$ & Pares generalizados. & N\textperiodcentered{}m\\
$m_i,\vct r_G$ & Masa de un cuerpo y centro de masa global. & kg, m\\
$\mat I_G,\mat M$ & Tensor de inercia central y matriz de inercia articular. & Según coordenadas\\
$K,U,\mathcal L$ & Energía cinética, potencial y lagrangiano. & J\\
$\mat\Sigma$ & Matriz de covarianza. & Producto de unidades\\
$k_x,k_y,k_\psi$ & Factores direccionales de odometría. & Adimensional\\
$i_r,\eta$ & Relación de reducción y eficiencia. & Adimensional\\
$g,r_p$ & Apertura de pinza y radio primitivo del piñón. & m\\
$g_0$ & Aceleración gravitatoria de referencia. & m/s$^2$\\
$R_m,L_m,K_e,K_t$ & Parámetros eléctricos y electromecánicos del motor. & Según definición\\
\bottomrule
\end{longtable}

\section*{Parámetros nominales usados}
\addcontentsline{toc}{section}{Parámetros nominales usados}
\begin{table}[!htbp]\centering
\begin{tabularx}{\linewidth}{@{}YP{50mm}@{}}\toprule
\textbf{Parámetro} & \textbf{Valor}\\\midrule
Diámetro y radio de rueda & $D=0{,}080\unit m$; $r=0{,}040\unit m$\\
Suma de semidistancias & $L=0{,}170\unit m$\\
Cuentas por vuelta, captura x1 & $N_i=370$\\
Altura equivalente de hombro & $h=0{,}209\unit m$\\
Eslabón proximal & $L_1=0{,}070\unit m$\\
Distancia equivalente a muñeca & $L_2=0{,}115\unit m$\\
Offset distal de herramienta & $d=0{,}1435\unit m$\\
Factores avance / lateral / giro & $1{,}000$ / $0{,}775$ / $1{,}028$\\
Período nominal de control & $10\unit{ms}$\\
Ganancias digitales iniciales & $K_p=128$, $K_i=8$, $K_d=0$; Q8\\\bottomrule
\end{tabularx}
\end{table}
Los valores de fábrica no sustituyen los guardados en el equipo. Las longitudes son nominales y las propiedades dinámicas y de contacto requieren caracterización.

\section*{Conversiones}
\addcontentsline{toc}{section}{Conversiones}
\begin{align}
1\unit{mm}&=10^{-3}\unit m, & 1^{\circ}&=\frac{\pi}{180}\unit{rad},\\
1\unit{rpm}&=\frac{2\pi}{60}\unit{rad/s}, &1\unit{kgf\,cm}&=0{,}0980665\unit{N\,m}.
\end{align}

\clearpage
\section*{Filminas de presentación del proyecto}
\addcontentsline{toc}{section}{Filminas de presentación del proyecto}
La presentación del proyecto se conserva en el repositorio \texttt{ProyectoRobotica2}, dentro de la carpeta \texttt{presentacion}. El archivo de las filminas es \texttt{AutoFinalRobo2\_Presentacion.pptx}.

\medskip
\noindent\href{https://github.com/kobachajchir/ProyectoRobotica2/blob/codex/monorepo-cad/presentacion/AutoFinalRobo2_Presentacion.pptx}{\textbf{Abrir las filminas en el repositorio de GitHub}}.

\medskip
\noindent\href{run:../../../../presentacion/AutoFinalRobo2_Presentacion.pptx}{\textbf{Abrir la copia local de las filminas}} desde este informe guardado en su carpeta de salida dentro del proyecto.

\medskip
Las filminas acompañan la exposición oral y reúnen arquitectura, modelos, conexionado, control, cinemática, odometría, esfuerzos e interfaces. La carpeta también contiene las fichas temáticas de apoyo y el guion de la exposición.

\clearpage
\begin{thebibliography}{9}
\addcontentsline{toc}{chapter}{Bibliografía}
\bibitem{lynch} Lynch, K. M. y Park, F. C. \emph{Modern Robotics: Mechanics, Planning, and Control}. Cambridge University Press, 2017. Material oficial de los autores: \url{https://hades.mech.northwestern.edu/index.php/Modern_Robotics}.
\bibitem{omni} Lynch, K. M. y Park, F. C. \emph{Omnidirectional Wheeled Mobile Robots}, sección 13.2 del material complementario de Modern Robotics. \url{https://modernrobotics.northwestern.edu/nu-gm-book-resource/13-2-omnidirectional-wheeled-mobile-robots-part-1-of-2/}.
\bibitem{mobile} Lynch, K. M. y Park, F. C. \emph{Mobile Manipulation}, sección 13.5 del material complementario de Modern Robotics. \url{https://modernrobotics.northwestern.edu/nu-gm-book-resource/13-5-mobile-manipulation/}.
\bibitem{servo} FEETECH. \emph{STS3250 Product Specification}. Ficha de la variante de servo utilizada para la selección nominal. \url{https://www.feetechrc.com/Data/feetechrc/upload/file/20240120/6384135881578380868917773.pdf}.
\end{thebibliography}
\endinput
